\section{Certifying an Alternative Matching}
\label{sec:prediction}

Fix a trace $\trace$, an alternative $\Malt$ with changed receives $\Changed$, and a
threshold $\Bt$. We first cut $\trace$ down to the part that every consistent program
repeats under $\Malt$ (the horizon), then decide $\Malt$ over that part with a
constraint graph.

\subsection{The horizon}
\label{sec:horizon}

\begin{definition}[Horizon]
The \emph{horizon} $\Hpre=\Hpre(\trace,\Malt,\Bt)$ is the greatest prefix of $\trace$
closed under the following rules. A record removed from a rank removes every later
record of that rank.
\begin{enumerate}[leftmargin=*,label=(H\arabic*)]
\item\label{H:changed} For $r\in\Changed$, the records of $\mathit{rank}(r)$ after the
  completion of $r$ are removed. Their results differ from $\trace$, so $\trace$ does not
  determine them.
\item\label{H:recv} A receive whose $\Malt$-send is not issued in $\Hpre$, or that has no
  $\Malt$-send, does not complete: the record completing it is removed.
\item\label{H:send} A send that blocks at $\Bt$ and whose $\Malt$-receive is not posted
  in $\Hpre$ does not complete. The record completing it is removed, but the send stays
  \emph{issued}, so a probe can still see it.
\item\label{H:coll} A collective whose shape requires a member that does not issue it in
  $\Hpre$ does not complete on the members that wait for it.
\item\label{H:probe} A probe whose observed send is not issued in $\Hpre$ is removed.
\item\label{H:pending} A receive $r'$ posted in $\Hpre$ whose $\Malt$-send is not issued in
  $\Hpre$ stays pending. Any later receive or probe of its rank that would take or report
  a message $r'$ accepts is removed, by \ref{E:order}.
\item\label{H:ready} A ready send whose $\Malt$-receive is not posted before it in
  $\Hpre$ is removed.
\end{enumerate}
The \emph{judged} receives $\judged$ are the receives posted in $\Hpre$ whose
$\Malt$-send is issued in $\Hpre$, together with the receives posted in $\Hpre$ that
$\Malt$ leaves unmatched. $\Malt$ \emph{survives} if every $r\in\Changed$ is judged.
\end{definition}

The horizon is computed as a fixpoint over per-rank cut indices. Rule~\ref{H:changed} is
the only rule that refers to $\Changed$. It encodes the one fact the trace cannot
supply: what a rank does after a call returns a result the trace did not record. Every
other rule removes an operation that could not complete if everything kept stayed as
recorded. A nonblocking changed receive is judged once it is posted and its new message
is issued, because the match is decided then. Its wait may lie beyond $\Hpre$.

\subsection{The constraint graph}
\label{sec:graph}

\begin{definition}[Constraint graph]
$\Graph=\Graph(\trace,\Malt,\Bt,\Hpre)$ has a node $\issue(e)$ for every operation issued
in $\Hpre$, a node $\compl(e)$ for every completion in $\Hpre$, and a node $\mnode(r,s)$
for every judged receive $r$ with $\Malt(r)=s$. Its edges are:
\begin{enumerate}[leftmargin=*,label=(G\arabic*)]
\item\label{G:po} program order on each rank, $\issue(e)\to\compl(e)\to\issue(e')$ for
  consecutive operations $e,e'$ (a nonblocking operation completes at its wait);
\item\label{G:match} for each match node, $\issue(r)\to\mnode(r,s)$,
  $\issue(s)\to\mnode(r,s)$, $\mnode(r,s)\to\compl(r)$, and $\mnode(r,s)\to\compl(s)$ if
  $\blocks_\Bt(s)$;
\item\label{G:fifo} for sends $s'$ before $s$ on one channel with $r\accepts s'$ and
  $\Malt(r)=s$: $\mnode(\Malt^{-1}(s'),s')\to\mnode(r,s)$;
\item\label{G:order} for receives $r'$ before $r$ on one rank with $r'\accepts s$ and
  $\Malt(r)=s$: $\mnode(r',\Malt(r'))\to\mnode(r,s)$;
\item\label{G:ready} for a ready send $s$ with $\Malt^{-1}(s)=r$: $\issue(r)\to\issue(s)$;
\item\label{G:coll} for each set of matching collectives, edges from issues to completions
  as their shape requires (through one hub node per set, so $O(p)$ edges);
\item\label{G:probe} for a probe that reports send $s$: $\issue(s)\to\compl(\mathit{probe})$, and
  edges into $\compl(\mathit{probe})$ from the matches \ref{G:fifo} and \ref{G:order} would
  require of a receive posted at the probe, since a probe reports what such a receive would
  take.
\end{enumerate}
A collective of a root shape without a recorded root synchronises all members.
A sender-side or receiver-side obligation whose earlier operation is issued in $\Hpre$
but unmatched under $\Malt$ makes the alternative infeasible. Rule~\ref{H:pending} has
already removed the receives that such a pending competitor would block.
\end{definition}

\begin{definition}[Feasibility]
$\Malt$ is \emph{feasible at $\Bt$} if it survives in $\Hpre(\trace,\Malt,\Bt)$, every
sender-side and receiver-side obligation over $\Hpre$ is met, and
$\Graph(\trace,\Malt,\Bt,\Hpre)$ is acyclic.
\end{definition}

Match nodes are what make the order rules exact. In MPI, the constraint that an earlier
receive be matched first binds the \emph{match} of the later receive. It does not bind
that receive's completion, which for a nonblocking receive lies at a wait much later, nor
the send's completion, which for a rendezvous send follows the match. Without match nodes,
\ref{G:order} has to be attached to some completion. Attached to $\compl(r)$, it misses
the cycle in which a rendezvous send must be matched before its sender can produce the
message the earlier receive was waiting for.
An earlier version of \tool did exactly that, and it accepted two shapes that MPI cannot
produce: a pending nonblocking wildcard overtaken by a later receive, and a rendezvous send
matched after its own completion. Both are now regression tests.

The edges of \ref{G:fifo} and \ref{G:order} need not be quadratic. The receives that accept
$s$ fall into four envelope classes (source $\mathit{rank}(s)$ or \ANY, tag $\mathit{tag}(s)$
or \ANY). Within a class, each receive accepts whatever a later one accepts. An edge from the
latest earlier member of each class therefore reaches the rest of the class through the
members' own edges. The same holds for the latest earlier send of each tag on $s$'s channel.

\begin{theorem}[Soundness]
\label{thm:sound}
If $\Malt$ is feasible at $\Bt$ with horizon $\Hpre$, then every program
$\prog\in\Cons(\trace,\Hpre)$ has an execution at $\Bt$ that issues every operation of
$\Hpre$, completes every operation whose completion lies in $\Hpre$, matches every judged
receive $r$ to $\Malt(r)$, and leaves the judged receives that $\Malt$ leaves unmatched
unmatched.
\end{theorem}

\noindent\emph{Proof idea.} Take a topological order of $\Graph$ and read it as a
schedule. Rule \ref{E:po} holds because each rank's results in $\Hpre$ equal those of
$\trace$, except at a changed receive. Each changed receive is the last record of its
rank in $\Hpre$ by \ref{H:changed}, so consistency makes $\prog$ issue exactly the
operations of $\Hpre$. When a match node is reached, \ref{G:fifo} and \ref{G:order}
guarantee that every earlier competitor on either side has been matched. Rule
\ref{H:pending} guarantees that no pending receive of the rank accepts the message, so
MPI's matching rule has no choice but $\Malt$. Rendezvous completions follow their
matches by \ref{G:match}, collectives follow their shapes, and probes see issued,
unmatched messages. \Cref{app:proofs} gives the full proof.
\pending{Proof being rewritten around match nodes.}

\subsection{Buffering profiles}
\label{sec:buffering}

Feasibility is a function of $\Bt$, and the thresholds that matter are $0$, the distinct
message sizes of $\trace$, and $\infty$. Let $\Bobs$ be the least threshold at which
$\Mobs$ itself is feasible. If $\Bobs>0$, the recorded run deadlocks on every
implementation that sends fewer than $\Bobs$ bytes eagerly, and \tool reports it as
\cls{2-buffering}.

\begin{lemma}
\label{lem:obsmono}
If $\Mobs$ is feasible at $\Bt$, it is feasible at every $\Bt'>\Bt$.
\end{lemma}
\noindent The horizon of $\Mobs$ is the whole trace. Raising the threshold only removes
the edges $\mnode(r,s)\to\compl(s)$ of \ref{G:match}.

\begin{theorem}[Monotonicity]
\label{thm:mono}
Let $\Malt$ be a steal and $\Bobs\le\Bt\le\Bt'$. If $\Malt$ is feasible at $\Bt$, it is
feasible at $\Bt'$, and $\Hpre(\trace,\Malt,\Bt)\subseteq\Hpre(\trace,\Malt,\Bt')$.
\end{theorem}

The hypothesis $\Bt\ge\Bobs$ is necessary. Below $\Bobs$, raising the threshold can bring
back into the horizon a head-to-head exchange that is cut at the lower threshold. That
exchange then closes a cycle until the threshold exceeds its messages, so feasibility can
switch from true to false and back. \Cref{thm:mono} makes bisection over
$[\Bobs,\infty]$ exact. \tool profiles each steal at $\Bobs$ and $\infty$ and bisects
between them. Judging every breakpoint is the differential test.
Swaps re-match two receives, which \cref{thm:mono} does not cover, so \tool profiles them at
every breakpoint. \pending{Proof: repair the step for a nonblocking changed receive, whose
judged receives can grow with $\Bt$.}

\begin{proposition}[Deadlocks are not monotone in $\Bt$]
\label{prop:nonmono}
The set of thresholds at which a predicted deadlock occurs is neither upward- nor
downward-closed in general.
\end{proposition}
\noindent Two witnesses suffice. A head-to-head exchange of 400-byte messages deadlocks
exactly for $\Bt<400$. In MBI's loop races, a steal needs the stolen message buffered,
so the resulting deadlock exists only for $\Bt\ge 4$. A deadlock prediction therefore
names its thresholds.

\subsection{Completeness for single changes}
\label{sec:complete}

Soundness alone is cheap: an analysis that accepts nothing is sound. Following the
optimality results for shared-memory race prediction~\citep{mathur21syncp}, we
characterise exactly which single changes the check accepts.

\begin{definition}[Matching-preserving reordering]
A \emph{matching-preserving reordering} of $\trace$ for a steal $r\mapsto s$ at $\Bt$ is an
execution at $\Bt$ of the operations of a prefix of $\trace$ that matches $r$ to $s$ and
matches every other receive it completes as $\Mobs$ does.
\end{definition}

\begin{theorem}[Completeness]
\label{thm:complete}
Let $\Bt\ge\Bobs$ and let the collectives of $\trace$ have known shapes. The check accepts
a steal $r\mapsto s$ at $\Bt$ if and only if a matching-preserving reordering for
$r\mapsto s$ at $\Bt$ exists.
\end{theorem}

\noindent\emph{Proof idea.} ``Only if'' is \cref{thm:sound} applied to $\SP(\trace)$.
For ``if'', every edge of $\Graph$ orders two events that every such reordering must
order the same way, and every rule of the horizon removes an operation that no such
reordering completes. A cycle would therefore lie inside the events the reordering must
contain.
\pending{T6 holds only after the conservative steps are removed or bounded: the global
cycle gate of the fast check, rule~\ref{H:probe}, freed receives counted as competitors,
and the order of an \mpicall{Mrecv} taken from its own posting. Proof to be written.}

\subsection{Cost}
\label{sec:cost}

\begin{theorem}[Cost]
\label{thm:cost}
Let $\trace$ have $n$ records, $w$ wildcard receives of which $w_q$ on rank $q$, at most $T$
tags on any channel read by an \ANY-tag receive, and $k$ distinct message sizes. Then
\begin{enumerate}[leftmargin=*,label=(\roman*)]
\item there are at most $w\cdot p$ steals and at most $\sum_q \binom{w_q}{2}$ swaps;
\item one alternative is decided at one threshold in time $O((n+|\Malt|\,T)\log n)$;
\item a steal is profiled with $O(\log k)$ checks above $\Bobs$ (\cref{thm:mono}), a swap with
  $k+2$.
\end{enumerate}
\end{theorem}

\noindent The bound on steals is the non-overtaking kernel: only the earliest unconsumed
message of each source can be next. The horizon is computed with a worklist. Each record
loses its issue at most once, because cuts only decrease, and each rule is re-evaluated only
when an input of its condition changes. Rule~\ref{H:pending} costs one bisection in an index
of the messages each receive takes. The graph has $O(n+|\Malt|\,T)$ edges thanks to the
class-indexed order edges, and cycle detection is linear. The check is therefore
quasi-linear per alternative, for any number of changed receives. What grows with the number
of changed receives is the number of alternatives. Choosing which receives to change is
NP-complete in general~\citep{forejt17mopper}, and placing that problem between polynomial
and hard cases parameterised by $w$~\citep{shi26popl,mathur20lics} is future work.
\pending{Tree clocks for the fast check's reachability; measured cost per alternative after
the rerun.}
